\documentclass[preprint,10pt,5p]{elsarticle}
\usepackage[utf8]{inputenc}
\usepackage[T1]{fontenc}
\usepackage[brazil]{babel}
\usepackage{graphicx}
\usepackage{amssymb}
\usepackage{amsmath}
\usepackage{hyperref}
\usepackage{float}
\usepackage{url}
\emergencystretch=3em
\journal{Universidade Católica Dom Bosco}

% Substitua o endereço abaixo por um link público, se houver.
\newcommand{\linkgraficos}{https://EXEMPLO-INSIRA-O-LINK-AQUI}

% Mostra um aviso em vez de interromper a compilação se a imagem não estiver enviada.
\newcommand{\imagemsegura}[2]{%
\IfFileExists{#1}{\includegraphics[width=#2\linewidth]{#1}}{\fbox{\parbox[c][3cm][c]{0.85\linewidth}{\centering Imagem pendente: \texttt{\detokenize{#1}}}}}%
}

\begin{document}
\begin{frontmatter}
\title{Análise Histórica de Indicadores da Educação Pública em Mato Grosso do Sul}
\author[ucdb]{Luis Felipe Andrade}
\ead{luis.silva170103@gmail.com}
\address[ucdb]{Universidade Católica Dom Bosco, Campo Grande, Mato Grosso do Sul, Brasil}

\begin{abstract}
Este relato apresenta a experiência de organização e visualização de indicadores da educação pública de Mato Grosso do Sul a partir de um projeto acadêmico desenvolvido anteriormente. Foram utilizados arquivos CSV, banco de dados PostgreSQL e uma aplicação em Python com Streamlit e Plotly para consultar e representar informações sobre matrículas e aprovação escolar. O material disponível contempla registros apresentados para o período de 2019 a 2026 e permite examinar a evolução das matrículas, comparar Campo Grande e Dourados e identificar municípios com maiores e menores taxas de aprovação. O artigo descreve o fluxo técnico, apresenta registros visuais do protótipo e discute suas possibilidades e limitações para a análise educacional. Os registros visuais não substituem uma nova execução ou validação do sistema.
\end{abstract}
\begin{keyword}
Educação pública \sep PostgreSQL \sep Streamlit \sep Visualização de dados \sep Indicadores educacionais
\end{keyword}
\end{frontmatter}

\section{Introdução}
A análise de dados educacionais contribui para a compreensão de diferenças entre localidades, variações no número de matrículas e indicadores de desempenho escolar. No contexto da educação pública de Mato Grosso do Sul, a organização de registros históricos pode apoiar a exploração de tendências e a formulação de perguntas sobre o funcionamento da rede de ensino.

A pandemia de COVID-19 alterou a rotina escolar e motivou o interesse pela observação de séries históricas antes, durante e após o período de restrições sanitárias. Entretanto, variações em indicadores não demonstram, isoladamente, uma relação causal com a pandemia. É necessário considerar a cobertura, a origem e a comparabilidade dos registros.

Este relato documenta uma experiência anterior de integração de dados CSV, PostgreSQL e ferramentas de visualização em Python. O objetivo é apresentar a organização técnica do projeto e os resultados exibidos em capturas do dashboard, reconhecendo as limitações do material recuperado.

\section{Trabalhos Relacionados}
A análise de dados educacionais pode combinar armazenamento relacional, consultas estruturadas e visualizações para facilitar a interpretação de informações de escolas e municípios. Bancos de dados relacionais permitem organizar registros e executar consultas, enquanto dashboards tornam mais acessível a exploração de indicadores por meio de filtros e gráficos.

A experiência relatada utiliza essas abordagens de forma integrada. O PostgreSQL foi empregado para armazenamento e consulta, enquanto Python, Streamlit e Plotly foram utilizados na apresentação dos indicadores. Esta seção descreve a relação conceitual entre as tecnologias empregadas, sem pretender constituir uma revisão sistemática da literatura.

\section{Materiais e Métodos}
\subsection{Dados e preparação}
O projeto anterior utilizou arquivos CSV contendo registros educacionais relacionados a ano, município, escola, matrículas e aprovação. Os arquivos foram organizados para importação ao banco de dados. A documentação preservada não permite detalhar todas as regras de limpeza nem verificar a completude da série histórica.

\subsection{Armazenamento no PostgreSQL}
Após a preparação, os registros foram importados para o PostgreSQL. A Figura~\ref{fig:banco} apresenta uma consulta à tabela \texttt{public.indicadores\_educacao}, com campos como identificador, ano, município e escola. A estrutura possibilitou consultar e agrupar informações para a etapa de visualização.

\begin{figure}[!htbp]
\centering
\imagemsegura{Postgree.png}{0.94}
\caption{Registro visual da consulta à tabela de indicadores educacionais no PostgreSQL, proveniente do projeto anterior.}
\label{fig:banco}
\end{figure}

\subsection{Aplicação e visualização}
O fluxo técnico adotado no projeto anterior foi: arquivos CSV, importação ao PostgreSQL, consulta e processamento com Python e apresentação no Streamlit com gráficos Plotly. O dashboard exibido nos registros possui filtros por ano e modalidades de análise individual, comparação e rankings. As capturas de tela são evidências da interface desenvolvida anteriormente, não uma comprovação de execução atual.

\section{Resultados}
\subsection{Evolução histórica das matrículas}
A Figura~\ref{fig:evolucao} apresenta a evolução das matrículas na série exibida pelo dashboard, de 2019 a 2026. Observa-se uma redução entre os pontos de 2019 e 2020 e uma nova diferença no valor apresentado para 2026. A figura não exibe registros para todos os anos intermediários; por isso, não é adequado afirmar que houve queda contínua em cada ano do período.

\begin{figure}[!htbp]
\centering
\imagemsegura{grafico1.png}{0.94}
\caption{Evolução histórica das matrículas exibida no dashboard do projeto anterior.}
\label{fig:evolucao}
\end{figure}

\subsection{Comparação entre municípios}
A Figura~\ref{fig:comparacao} apresenta os indicadores de Campo Grande e Dourados para o ano selecionado de 2019. Na captura, Campo Grande registra 70.087 matrículas e 45.045 aprovados, com taxa de aprovação exibida de 64,27\%. Dourados apresenta 24.396 matrículas e 16.077 aprovados, com taxa de 65,90\%. A comparação ilustra a diferença entre quantidades absolutas e indicadores percentuais.

\begin{figure}[!htbp]
\centering
\imagemsegura{grafico2.png}{0.94}
\caption{Comparação de matrículas, aprovados e taxas de aprovação entre Campo Grande e Dourados, com filtro de 2019.}
\label{fig:comparacao}
\end{figure}

\subsection{Rankings de aprovação}
As Figuras~\ref{fig:maiores} e~\ref{fig:menores} apresentam rankings de municípios com maiores e menores taxas de aprovação no ano selecionado de 2019. Os gráficos demonstram a funcionalidade de classificação do dashboard e evidenciam diferenças entre os indicadores municipais. A interpretação desses rankings depende da definição do indicador, da quantidade de matrículas e da qualidade dos registros de origem.

\begin{figure}[!htbp]
\centering
\imagemsegura{grafico3.png}{0.94}
\caption{Municípios com maiores taxas de aprovação, segundo o dashboard anterior, com filtro de 2019.}
\label{fig:maiores}
\end{figure}

\begin{figure}[!htbp]
\centering
\imagemsegura{grafico4.png}{0.94}
\caption{Municípios com menores taxas de aprovação, segundo o dashboard anterior, com filtro de 2019.}
\label{fig:menores}
\end{figure}

% Se tiver um link público, retire o % da linha seguinte e substitua o endereço no início do arquivo.
% \noindent\textbf{Visualização complementar:} \href{\linkgraficos}{acessar os registros visuais do projeto}.\par

\section{Discussão}
O principal resultado da experiência foi a integração entre armazenamento relacional e visualização interativa. A possibilidade de alternar entre análise histórica, comparação de municípios e rankings oferece diferentes perspectivas sobre o mesmo conjunto de dados.

A série de matrículas sugere variações relevantes ao longo do período representado, mas não permite atribuí-las exclusivamente à pandemia. Além disso, os registros recuperados não demonstram cobertura completa de todos os anos e não incluem uma validação independente dos valores. Na comparação entre Campo Grande e Dourados, a maior quantidade absoluta de matrículas em Campo Grande não corresponde à maior taxa de aprovação exibida, o que evidencia a importância de observar indicadores relativos.

Como limitações, destaca-se a indisponibilidade, neste momento, do código-fonte completo e dos arquivos originais de dados para uma nova execução e auditoria. Também não foi documentada nesta versão uma implementação específica de orientação a objetos ou de estruturas de dados. Assim, o texto registra a experiência anterior e seus resultados visuais, sem apresentar o material como um protótipo reconstruído ou validado neste semestre.

\section{Conclusão}
A experiência demonstrou a aplicação conjunta de PostgreSQL, Python, Streamlit e Plotly na organização e visualização de indicadores da educação pública de Mato Grosso do Sul. Os registros disponíveis mostram funcionalidades de análise histórica, comparação entre municípios e classificação por taxas de aprovação.

O trabalho também evidencia a necessidade de preservar código-fonte, bases de dados e documentação de processamento para garantir reprodutibilidade. Como continuidade, seria importante recuperar os arquivos originais, verificar a cobertura temporal, documentar os critérios de cálculo e validar os indicadores antes de produzir conclusões sobre os efeitos da pandemia.

\end{document}
